\section{Grothendieck's questions and the results}
The distinction between a function field and an affine model is decisive
for the questions considered here. In cote 31 of the Montpellier archive,
Grothendieck proposes bounds for the cohomological dimension of an affine
scheme in terms of its function field and residue fields. On archive
sheet 84, headed \emph{Conjectures sur la dimension cohomologique des
sch\'emas affines}, he writes, for an integral affine scheme with
$\ell$ invertible,
\[\cd_\ell(X)\leq\cd_\ell(k(X)).\]
The same sheet proposes the residual-characteristic affine-open bound
answered below \cite[sheet 84]{Grothendieck31}. We use $\cd_\ell(X)$ for the least integer $d$ such
that $H^q_{\mathrm{et}}(X,\mathcal F)=0$ for every $\ell$-primary torsion
sheaf $\mathcal F$ and every $q>d$. This convention includes arbitrary
torsion sheaves, not only constant coefficients.

\begin{theorem}\label{thm:arithmetic}
There is an excellent Dedekind domain $R$, of characteristic zero with
$2\in R^\times$, such that its fraction field $K$ is algebraic over $\Q$,
\[
 \cd_2(K)=1,\qquad
 \cd_2(U)=2\quad\text{for every nonempty affine open }U\subset\Spec R.
\]
For every nonempty open $U\subset\Spec R$, the vector space
$\Pic(U)/2$ over $\F_2$ has infinite dimension. No residue field of $R$
is orderable, and
\[
 \sup_{x\in U}\{\cd_2\kappa(x)+\dim\OO_{U,x}\}=1.
\]
\end{theorem}

Thus no shrinking removes the gap between the cohomological dimensions
of this regular arithmetic curve and its generic point. Invertibility
of two, excellence and the absence of real residue fields all hold. The
mechanism is a line-bundle obstruction: Kummer theory places
$\Pic(U)/2$ inside $H^2(U,\mu_2)$.

The construction is based on the reciprocity, ray-class existence and
transfer theorems of class field theory, in the precise forms recalled
by Milne \cite{MilneCFT}. It preserves old ideal classes modulo two
while at least doubling their dimension. Simultaneously, eventual
inertness at every old prime makes the limit ring noetherian and kills
the two-primary Brauer groups of all finite extensions of its function
field. These two effects must be arranged together. Claborn proves that
every abelian group is the class group of a Dedekind domain
\cite[Theorem 7, p.~222]{Claborn}. That realization theorem does not assert
the simultaneous algebraic-number-field, residue-field and cohomological
conditions imposed here.
S.~A.~Popescu already constructs an eventually inert quadratic tower
whose integral closure of $\Z[1/2]$ is Dedekind
\cite[\S5, pp.~511--512]{Popescu2022}. His construction controls
the odd primes. The quadratic step below adds control at dyadic primes
and injectivity of old ideal classes modulo two, together with a doubling
of their dimension. These additions supply the two cohomological features.

\begin{theorem}\label{thm:wild}
Let $k=\overline{\F}_2$, let $W(k)$ be its Witt ring, and set
\[
 A=W(k)[[t]],\qquad U=\Spec A[1/t].
\]
Then $A$ is a complete excellent regular local ring of dimension two,
but $H^3(U,\F_2)\ne0$ and $\cd_2(U)=3$.
\end{theorem}

Here the omitted divisor has characteristic-zero generic point, while
the open scheme still meets characteristic two. The nonzero class is
the boundary of $\{1+2/t,t\}$ in a henselian height-one field. Kato's
ramification filtration, as computed by Suzuki
\cite[Propositions 6.1.2, 8.2.1 and 8.2.3]{Suzuki}, separates this
symbol from classes extending across the other divisor. This example
answers the residual-characteristic affine-open question in the same
notes. It does not concern the prime-to-residue-characteristic affine
vanishing theorem.

In cote 29, archive sheet 70, Grothendieck states the Picard assertion
as a theorem:
\begin{quote}
\guillemotleft\ Alors l'homomorphisme $\Pic^\tau_Y\to\Pic^\tau_X$
est un isomorphisme\guillemotright.
\end{quote}
The argument continues on sheet 71 \cite[sheets 70--71]{Grothendieck29}.
For a proper
variety $Z$, $\Pic^\tau_{Z/k}$ denotes the union of the components of
the Picard scheme whose N\'eron--Severi classes are torsion.

\begin{theorem}\label{thm:birational}
Let $f:X\to Y$ be a proper birational morphism of integral noetherian
schemes of characteristic $p>0$. If $X$ is normal and $Y$ is regular,
then $R^1f_*\OO_X=0$.

If, in addition, $X$ and $Y$ are proper varieties over an algebraically
closed field $k$, then pullback induces an isomorphism of group schemes
\[
 \Pic^\tau_{Y/k}\xrightarrow{\sim}\Pic^\tau_{X/k}.
\]
\end{theorem}

The coherent vanishing is already Ishii--Yoshida's theorem
\cite[Corollary 1.4, p.~3]{IshiiYoshida}, in the stronger form allowing
every characteristic. The Picard assertion is also a consequence of known
work: Li--Wang prove invariance of the Nori fundamental group scheme for
a normal proper source and smooth target over a perfect field
\cite[Corollary 4.6, p.~7]{LiWang}. Their statement includes the finite
commutative torsors that determine the torsion Picard scheme. We explain
this last deduction, using the published purity theorem directly.
Grothendieck's argument on sheet 71 reduces the
Picard problem to its infinitesimal structure \cite{Grothendieck29}. Points and tangent
spaces alone cannot finish that argument: the inclusion
$\alpha_p\subset\alpha_{p^2}$ is bijective on geometric points and
on tangent spaces, but its two coordinate algebras have dimensions
$p$ and $p^2$. Purity for finite flat torsors controls every finite
subgroup, including this higher infinitesimal structure. The general
purity theorem is due to Moret--Bailly; Marrama gives a published
detailed proof \cite[Theorem 6, pp.~433--434]{Marrama}.
The quoted vanishing theorem and the finite-torsor argument both
control infinitesimal information, but the Picard assertion requires
all finite subgroup schemes, not merely a tangent-space computation.

Recent work on Picard groups also studies faithfully flat changes of the
coefficient ring. The Hexagon note contributed by Mathew and produced by
ChatGPT~6 Pro proves that every invertible $R[T]$-module becomes free
after a faithfully flat extension of $R$
\cite[Theorem 1.1, p.~1]{HexagonPicard}, using Bhatt's graded results.
That theorem concerns polynomial extensions and arbitrary flat base
change; the Picard assertion here concerns a fixed proper birational map.
It is a contextual comparison and is not an input to our proofs.

Finally, for a normal proper connected variety $X$ over an algebraically
closed field, and a proper closed subset $Y$, let $Q(X,Y)$ be the
quotient of $\pi_1(X\setminus Y)$ by the closed normal subgroup
generated by wild inertia at the codimension-one components of $Y$.
This definition uses the given compactification.

\begin{theorem}\label{thm:boundary}
If $Y$ is contained in a quasi-projective open subset of $X$, then
$Q(X,Y)$ is topologically finitely generated. In particular,
$\pi_1(X\setminus\{x\})$ is topologically finitely generated when
$x$ is a closed point of codimension at least two. More generally,
if a proper connected variety $Z$ need not be normal and a closed
subset has codimension at least two on every irreducible component
and has one quasi-projective open neighborhood, the fundamental group
of its complement is finitely generated.

For normal $X$, the same conclusion for $Q(X,Y)$ holds whenever
there is a normal projective modification $X'\to X$ such that the
wild inertia at each exceptional divisorial component of the inverse
image of $Y$ has finite image in $Q(X,Y)$.
\end{theorem}

The point is the prescribed neighborhood in Raynaud's form of Chow's
lemma \cite[Chapter 4, \S4, Proposition 5, p.~30]{Raynaud}. All
exceptional divisors of the modification then lie in the open scheme,
where the covers are already \'etale. The proof therefore avoids
introducing uncontrolled boundary inertia. In the projective case,
the curve argument is already present in Drinfeld
\cite[Lemma C.2, p.~19]{Drinfeld}; Esnault--Kindler extend its
finite-field formulation to every field
\cite[Proposition 6.2, pp.~8--9]{EsnaultKindler}. Thus the projective
input is a known Lefschetz consequence. We give the deduction for an
arbitrary closed boundary, using the connectedness theorem of
Martinelli--Naranjo--Pirola \cite[Theorem 1.1, p.~31]{MNP} in the same
curve argument. The tame-curve theorem of SGA~1 supplies finitely many
generators. The neighborhood assertion then follows from the prescribed
Chow modification. The final criterion uses one finite cover that
removes the finitely many wild inertia images.
